\documentclass[11pt,a4paper]{article}
\usepackage[T1]{fontenc}
\usepackage[utf8]{inputenc}
\usepackage[french]{babel}
\usepackage{lmodern}
\usepackage{amsmath,amssymb,mathtools}
\usepackage{geometry}
\geometry{top=22mm,bottom=22mm,left=23mm,right=23mm,headheight=15pt,headsep=6mm}
\usepackage{microtype,enumitem,booktabs,array,fancyhdr,lastpage,needspace}
\usepackage[hidelinks]{hyperref}
\setlength{\parindent}{0pt}
\setlength{\parskip}{6pt}
\setlist[enumerate]{leftmargin=*,label=\textbf{\alph*)},itemsep=8pt,topsep=4pt}
\pagestyle{fancy}\fancyhf{}
\renewcommand{\headrulewidth}{0.3pt}
\fancyfoot[C]{\small\thepage\,/\,\pageref{LastPage}}
\fancyfoot[L]{\scriptsize Version de travail -- 6 septembre 2026}
\setlength{\emergencystretch}{2em}
\hypersetup{pdftitle={MAT0339 -- Série C, examen 3 -- sujet},pdfauthor={Document de travail pédagogique}}
\fancyhead[L]{\small MAT0339 -- Série C -- Examen 3}
\fancyhead[R]{\small SUJET À VALIDER}
\begin{document}
{\LARGE\bfseries Série C -- Examen 3}\par
{\large Optimisation, trigonométrie et synthèse}\par
\textbf{Épreuve étudiante -- version de travail}\par
Portée : chapitres 13 à 17; rappel cumulatif de 20 points. Total : 100 points. Durée proposée : 150 min.\par
\textbf{Nom :} \hrulefill\quad \textbf{Groupe :} \rule{22mm}{0.3pt}\par
\small Montrer les démarches. Conserver les valeurs exactes jusqu’à l’arrondi final; annoncer les unités. Convention provisoire : calculatrice scientifique non symbolique autorisée, sans notes. Les modalités doivent être confirmées par l’enseignant. Les pages suivantes offrent une page par question; une feuille supplémentaire peut être jointe.\normalsize\par
\vspace{3mm}\hrule\vspace{4mm}
{\large\bfseries Question 1 -- Optimisation à deux variables}\hfill\textbf{20 points}\par
\begin{enumerate}
\item \textbf{(5 points)} Un atelier prépare deux mélanges divisibles, en quantités \(x,y\ge0\). Les besoins par unité et les disponibilités sont :
\begin{center}\begin{tabular}{lrrr}\toprule
& Mélange X & Mélange Y & Disponible\\\midrule
Ressource 1 & 2 & 1 & 14\\
Ressource 2 & 1 & 3 & 18\\
Bénéfice (\$ par unité) & 40 & 50 & ---\\\bottomrule
\end{tabular}\end{center}
Écrire les contraintes et la fonction objectif à maximiser. Les quantités ne sont pas nécessairement entières.
\item \textbf{(5 points)} Tracer la région réalisable et calculer tous ses sommets.
\item \textbf{(6 points)} Évaluer l’objectif aux sommets et déterminer le maximum.
\item \textbf{(4 points)} Vérifier l’admissibilité du point optimal et préciser les ressources entièrement utilisées.
\end{enumerate}
\vspace{3mm}\textit{Démarche et réponse :}\par\vfill

\newpage
{\large\bfseries Question 2 -- Radians, cercle et secteur}\hfill\textbf{15 points}\par
\begin{enumerate}
\item \textbf{(3 points)} Convertir exactement l’angle \(\theta=330^\circ\) en radians.
\item \textbf{(4 points)} Donner les valeurs exactes de \(\sin\theta\) et \(\cos\theta\), en justifiant les signes par le quadrant.
\item \textbf{(4 points)} Un secteur circulaire a un rayon de 5 cm et un angle au centre \(\alpha=4\pi/5\) radians. Calculer la longueur de son arc.
\item \textbf{(4 points)} Calculer l’aire de ce même secteur. Conserver une valeur exacte.
\end{enumerate}
\vspace{3mm}\textit{Démarche et réponse :}\par\vfill

\newpage
{\large\bfseries Question 3 -- Résoudre un triangle}\hfill\textbf{15 points}\par
\begin{enumerate}
\item \textbf{(6 points)} Dans un triangle \(ABC\), les côtés \(a=5\) cm et \(b=8\) cm sont opposés aux angles \(A\) et \(B\), et l’angle compris vaut \(C=60^\circ\). Faire un schéma étiqueté, choisir une loi et déterminer exactement \(c\).
\item \textbf{(5 points)} Déterminer l’angle \(A\) au centième de degré. Donner une formule qui permet d’éviter une ambiguïté d’arcsinus.
\item \textbf{(4 points)} Calculer exactement l’aire du triangle et vérifier l’inégalité triangulaire pour le côté \(c\).
\end{enumerate}
\vspace{3mm}\textit{Démarche et réponse :}\par\vfill

\newpage
{\large\bfseries Question 4 -- Toutes les solutions, pas seulement l’angle principal}\hfill\textbf{15 points}\par
\begin{enumerate}
\item \textbf{(7 points)} Résoudre \(\tan x=1\) sur \([0,2\pi[\). Justifier que la liste est complète.
\item \textbf{(8 points)} Résoudre \(2\sin^2x-1=0\) sur \([0,2\pi[\). Donner des valeurs exactes.
\end{enumerate}
\vspace{3mm}\textit{Démarche et réponse :}\par\vfill

\newpage
{\large\bfseries Question 5 -- Construire puis utiliser une sinusoïde}\hfill\textbf{15 points}\par
\begin{enumerate}
\item \textbf{(9 points)} Une hauteur de marée varie sinusoïdalement entre 1 m et 7 m, avec une période de 12 heures. Au temps \(t=2\), la hauteur franchit la ligne médiane en montant. Construire une formule en sinus ou cosinus. Identifier amplitude, médiane et période, puis vérifier les données caractéristiques.
\item \textbf{(6 points)} Trouver tous les instants où la hauteur vaut 5,5 m pendant \([0,12[\) heures.
\end{enumerate}
\vspace{3mm}\textit{Démarche et réponse :}\par\vfill

\newpage
{\large\bfseries Question 6 -- Réinvestissement cumulatif}\hfill\textbf{20 points}\par
\begin{enumerate}
\item \textbf{(6 points)} Résoudre \((x-2)/(x+1)>0\) par les signes et préciser les bornes.
\item \textbf{(6 points)} Déterminer la réciproque de \(g(x)=5-x\), de \(\mathbb R\) vers \(\mathbb R\), et vérifier les deux compositions.
\item \textbf{(8 points)} On lance trois fois une pièce équilibrée, indépendamment. Calculer la probabilité d’obtenir exactement deux faces et justifier le facteur de dénombrement.
\end{enumerate}
\vspace{3mm}\textit{Démarche et réponse :}\par\vfill
\end{document}
